\chapter{Strings, Symbols and Text}
\label{ch:text}

\standfirst{Three classes that look alike and are not interchangeable, and one
of them carries formatting.}

\section{The three}

\begin{center}
\small
\begin{tabular}{@{}lp{0.62\textwidth}@{}}
\toprule
\ctp{String} & a mutable, byte-indexed collection of Characters\\
\ctp{Symbol} & an interned, immutable String. Two symbols spelled alike are
  the \emph{same object}\\
\ctp{Text} & a String plus runs of emphasis---bold, italic, underline. What
  the Browser and the editors actually manipulate\\
\bottomrule
\end{tabular}
\end{center}

\ct{Symbol} is a subclass of \ct{String}, so every String message works on a
symbol; the reverse is not true, because a symbol may not be modified.

\subsection{When to use which}

Use a \ct{Symbol} for anything that names something---a selector, a key in a
dictionary, a state in a state machine, a tag. Comparison is pointer
comparison and hashing is free. Use a \ct{String} for data: something a user
typed, something read from a file, something you built.

\begin{st}
#hello == #hello             "true  -- one object"
'hello' == 'hello'           "false -- two objects"
'hello' = 'hello'            "true"
'hello' asSymbol             "#hello"
#hello asString              "'hello'"
\end{st}

\section{String protocol}

A String is a collection, so everything in Chapter~\ref{ch:collections}
applies: \ct{do:}, \ct{select:}, \ct{collect:}, \ct{detect:}, \ct{size},
\ct{includes:}, \ct{,}, \ct{copyFrom:to:}, \ct{indexOf:}, \ct{reverse}.

On top of that:

\begin{center}
\small
\begin{tabular}{@{}ll@{}}
\toprule
\ctp{asUppercase} \ctp{asLowercase} & case\\
\ctp{asSymbol} \ctp{asNumber} \ctp{asInteger} & conversion\\
\ctp{asArray} \ctp{asBag} \ctp{asSet} & conversion\\
\ctp{reversed} \ctp{sorted} & \ctp{'hello' sorted} is \ctp{'ehllo'}\\
\ctp{copyReplaceAll:with:} & substitution\\
\ctp{indexOfSubCollection:startingAt:} & find a substring\\
\ctp{findString:startingAt:} & the same thing, String's own\\
\ctp{includesSubstring:} & \\
\ctp{trimBoth} & strip surrounding whitespace\\
\ctp{isEmpty} \ctp{notEmpty} & \\
\ctp{displayString} \ctp{printString} & the second adds quotes\\
\ctp{withCRs} & turn backslashes into carriage returns\\
\bottomrule
\end{tabular}
\end{center}

\begin{st}
'hello' asUppercase                        "'HELLO'"
'hello' reversed                           "'olleh'"
'hello' indexOf: $l                        "3"
'hello' copyFrom: 2 to: 4                  "'ell'"
'abcabc' copyReplaceAll: 'b' with: 'X'     "'aXcaXc'"
'hello' , ' ' , 'world'                    "'hello world'"
\end{st}

\begin{stnote}[Line endings]
Smalltalk-80 uses carriage return, \ct{Character cr}, as its line separator,
not line feed. \ct{withCRs} turns \ct{\textbackslash} into a carriage return,
which is how the 1983 sources write multi-line menu labels on one line.
\end{stnote}

\section{Characters}

\begin{st}
$a                          "a Character literal"
$a asciiValue               "97"
$a asInteger                "97"
Character value: 97         "$a"
$a asUppercase              "$A"
$a isVowel                  "true"
$a isLetter                 "true"
$5 digitValue               "5"
Character cr
Character tab
Character space
\end{st}

\begin{stwarn}
\ct|$a value| answers \ct|$a|, not 97: \ct{Character} implements
\ct{asciiValue}, and \ct{Object>>value}---added for Pharo
compatibility---answers the receiver. Use \ct{asciiValue} or \ct{asInteger}.
\end{stwarn}

\section{Building strings}

Concatenation with \ct{,} is fine for two or three pieces. For a loop, use a
stream (Chapter~\ref{ch:streams}):

\begin{st}
| s |
s := WriteStream on: (String new: 0).
1 to: 5 do: [:i | s print: i; nextPut: $ ].
s contents                        "'1 2 3 4 5 '"
\end{st}

\section{Text and emphasis}

\ct{Text} pairs a String with runs of \ct{TextAttribute}. The editors,
the Browser's code pane and \ct{Transcript} all deal in Text rather than
String, which is why some methods answer \ct{aText} where you expected a
string.

\begin{st}
'hello' asText
aText asString                   "back to plain characters"
\end{st}

You will meet \ct{Text} mostly by accident---as the thing a view hands you.
\ct{asString} is the usual response.

\section{Fonts, and why they matter here}

Text is drawn from a \ct{StrikeFont}: one wide one-bit bitmap of every glyph
side by side, plus an \ct{xTable} of where each one starts. The width of a
character is the difference between two adjacent entries, which is what makes
the face proportional.

The built-in face is Inter at 18 pixels with three rows of lead, rasterised
into a checked-in C file. The screen shows an antialiased version of it, but
the \emph{image} only ever sees one bit per pixel---Chapter~\ref{ch:graphics}
explains how both can be true.

\begin{stwarn}
Changing the face means rebuilding the image. \ct{TextList} fixes its line
grid from the font's height at class-initialisation time, and \ct{PopUpMenu}
composes its labels into a Form once, so an image built against one face
cannot be shown another. The VM says so at startup if they disagree.
\end{stwarn}
